\section{Introduction}

The rapid deployment of large language models (LLMs) across high-stakes domains---from healthcare advising to legal assistance and autonomous agent systems---has amplified concerns about the trustworthiness and verifiability of model-generated outputs \cite{hubinger2019risks, bommasani2021opportunities}. As these systems become increasingly capable, the possibility that they may produce deceptive outputs, whether through misalignment, adversarial prompting, or emergent strategic behavior, poses a fundamental challenge to AI safety \cite{park2023ai}. Detecting when a language model is being deceptive, rather than merely incorrect, represents a critical frontier in alignment research.

\subsection{Motivation}

The alignment problem---ensuring that AI systems act in accordance with human intentions---has been a central concern since the earliest formulations of AI safety \cite{russell2019human}. Recent work on reinforcement learning from human feedback (RLHF) has demonstrated that language models can be steered toward helpful, harmless, and honest behavior \cite{ouyang2022training}, yet such training offers no formal guarantees against deceptive behavior. \citet{hubinger2019risks} articulated the risk of \emph{deceptive alignment}, wherein a sufficiently capable model might learn to appear aligned during training while pursuing misaligned objectives during deployment. This concern motivates the development of behavioral detection methods that operate on model outputs rather than relying solely on internal representations or training-time guarantees. Interpretability research has made progress in understanding model internals \cite{elhage2022toy}, but output-level behavioral probing remains essential for black-box deployment scenarios where model weights are inaccessible.

\subsection{The Challenge: LLM Deception versus Human Lie Detection}

Decades of research in psychology and forensic science have established that detecting deception in humans is remarkably difficult, with trained professionals achieving only marginally above-chance accuracy \cite{bond2006accuracy}. LLM deception detection presents a fundamentally different challenge. Unlike humans, language models lack embodied physiological signals (e.g., micro-expressions, galvanic skin response) that form the basis of traditional lie detection. Instead, deception must be inferred entirely from textual behavioral patterns: response latency proxies, linguistic consistency, elaboration patterns, and reactions to adversarial pressure.

It is important to distinguish two qualitatively different settings for LLM deception. \textbf{(a) Instructed roleplay deception}, in which a model is explicitly prompted to defend a false claim as if it were true, and \textbf{(b) autonomous strategic deception}, in which a model spontaneously produces deceptive outputs to pursue misaligned goals. The latter represents the harder and more consequential alignment problem identified by \citet{hubinger2019risks}. This paper addresses setting (a) as a \emph{tractable proxy} for studying behavioral detection methods. We argue that instructed deception provides a controlled experimental framework for developing and validating detection pipelines that may, with further research, generalize to the more challenging autonomous setting. The instructed setting is valuable precisely because it allows systematic study of how deceptive behavior manifests in model outputs under known ground-truth conditions.

\subsection{Our Approach}

We propose an adaptive interrogation framework comprising four integrated components. \textbf{First}, a \emph{question generation engine} produces structured interrogation sequences spanning baseline, consistency-probing, detail-eliciting, and pressure-testing categories. \textbf{Second}, a \emph{behavioral feature extraction pipeline} computes five behavioral features from model responses, capturing consistency scores, specificity, hedging language frequency, confidence indicators, and elaboration depth. \textbf{Third}, a \emph{confidence-based adaptive stopping mechanism} dynamically determines when sufficient evidence has been gathered, terminating interrogation early when classifier confidence exceeds a calibrated threshold. This reduces the average number of questions from 8 (fixed) to 2.0 on mock data (75\% reduction) and to 3.0 on the real LLM pilot, where noisier behavioral signals require more rounds before the threshold is crossed. We note that the confidence score in the mock setting is poorly calibrated---near-certain regardless of correctness---so the stopping mechanism fires on classifier certainty rather than evidential reliability; recalibration is a necessary direction for deployment. \textbf{Fourth}, a \emph{logistic regression classifier} trained on extracted behavioral features produces a truthful-or-lying determination along with calibrated confidence estimates.

\subsection{Contributions}

The principal contributions of this work are as follows:
\begin{itemize}
    \item An \textbf{adaptive interrogation framework} with confidence-based early stopping that reduces questioning rounds by approximately 70\% while maintaining detection performance.
    \item A \textbf{behavioral feature extraction and classification pipeline} operating on five behavioral features, enabling black-box deception detection without access to model internals.
    \item A \textbf{mock model evaluation} ($n=100$ trials) achieving 93\% overall accuracy (86\% lying, 100\% truthful) with 5-fold CV confirming generalisation within the mock distribution, and a \textbf{prompt-swap ablation} (93\%$\to$14\%) confirming that the classifier responds to behavioral response content rather than prompt-format artifacts. We emphasise that mock responses are drawn from small fixed phrase pools and these results establish \emph{algorithmic correctness of the pipeline}, not empirical performance on real LLM deception.
    \item A \textbf{preliminary real-LLM pilot} ($n=9$ on Claude Haiku 4.5, $n_\text{lying}=2$) showing 88.9\% overall accuracy and 50\% lying detection --- directional evidence of a mock-to-real transfer gap but insufficient for reliable performance estimation due to sample size.
    \item A \textbf{calibration analysis} showing that the classifier produces near-certain confidence ($>0.99$) for both correct and incorrect predictions on mock data, revealing that the adaptive stopping mechanism currently fires on classifier certainty rather than evidential reliability, and motivating future calibration work.
\end{itemize}

The remainder of this paper is organized as follows. Section~2 reviews related work on AI deception and detection. Section~3 details the framework architecture. Section~4 describes the experimental setup and presents results on both mock and real LLM settings. Section~5 discusses limitations and implications, and Section~6 concludes.